\chapter{عیب‌یابی و رفع خطا}

با افزایش پیچیدگی اسکریپت‌ها، مدیریت و شناسایی خطاهای احتمالی اهمیتی دوچندان می‌یابد. در این فصل، به بررسی دسته‌بندی‌های رایج خطا در اسکریپت‌نویسی پوسته می‌پردازیم و تکنیک‌های کارآمدی را برای شناسایی (\en{Debugging}) و رفع این اختلالات مرور خواهیم کرد.

\section{خطاهای نحوی (Syntax Errors)}

یکی از رایج‌ترین دسته‌های خطا، «خطاهای نحوی» هستند که بر اثر اشتباه در نگارش ساختارهای دستوریِ پوسته پدید می‌آیند. مفسر پوسته به محض مواجهه با چنین خطاهایی، فرآیند اجرای اسکریپت را متوقف می‌کند. برای نمایش انواع متداول این خطاها، از اسکریپت نمونه‌ی زیر استفاده می‌کنیم:

\begin{latin}
\begin{lstlisting}
#!/bin/bash

# trouble: script to demonstrate common errors

number=1

if [ $number = 1 ]; then
    echo "Number is equal to 1."
else
    echo "Number is not equal to 1."
fi
\end{lstlisting}
\end{latin}

این اسکریپت در وضعیت فعلی، بدون هیچ‌گونه نقص فنی اجرا می‌شود:

\begin{latin}
\begin{lstlisting}
[me@linuxbox ~]$ trouble
Number is equal to 1.
\end{lstlisting}
\end{latin}

\subsection{خطاهای کاراکتر کوتیشن}

اکنون اسکریپت را ویرایش کرده و کاراکتر کوتیشن پایانی را در آرگومان دستور \cmd{echo} اول حذف می‌کنیم:

\begin{latin}
\begin{lstlisting}
#!/bin/bash

# trouble: script to demonstrate common errors

number=1

if [ $number = 1 ]; then
    echo "Number is equal to 1.
else
    echo "Number is not equal to 1."
fi
\end{lstlisting}
\end{latin}

خروجی مفسر پس از اجرا به این صورت خواهد بود:

\begin{latin}
\begin{lstlisting}
[me@linuxbox ~]$ trouble
/home/me/bin/trouble: line 10: unexpected EOF while looking for matching `"'
/home/me/bin/trouble: line 13: syntax error: unexpected end of file
\end{lstlisting}
\end{latin}

مشاهده می‌کنید که دو خطا تولید شده است. نکته‌ی قابل‌تأمل اینجاست که شماره‌ خطوط گزارش‌شده در پیام‌های خطا، دقیقاً به محل حذف کوتیشن اشاره ندارند، بلکه نقاطی بسیار جلوتر در برنامه را نشان می‌دهند. با تحلیل رفتار مفسر، علت این موضوع مشخص می‌شود: \en{Bash} پس از مواجهه با کوتیشنِ باز، تمام نویسه‌های بعدی را به عنوان بخشی از رشته در نظر می‌گیرد تا زمانی که به یک کوتیشنِ بسته برسد؛ در این مثال، مفسر اولین کوتیشنِ بسته را در دستور \cmd{echo} دوم می‌یابد. در نتیجه، ساختار دستور \cmd{if} از هم می‌پاشد، چرا که کلمه‌ی کلیدی \cmd{fi} اکنون در میان یک رشته‌ی متنیِ ناتمام محصور شده است.

در اسکریپت‌های طولانی و پیچیده، ردیابی این قبیل خطاها می‌تواند بسیار دشوار باشد. استفاده از ویرایشگرهای متنی که از قابلیت برجسته‌سازی نحو (\en{Syntax Highlighting}) پشتیبانی می‌کنند، کمک شایانی به رفع این مشکل می‌کند؛ زیرا رشته‌های متنی را با رنگی متفاوت از دستورات اصلی پوسته نمایش می‌دهند. چنانچه نسخه‌ی کامل ویرایشگر \cmd{vim} بر روی سیستم شما نصب است، می‌توانید با وارد کردن دستور زیر در محیط برنامه، این قابلیت را فعال کنید:

\begin{latin}
\begin{lstlisting}
:syntax on
\end{lstlisting}
\end{latin}

\subsection{خطاهای توکن‌های مفقود یا غیرمنتظره}

یکی دیگر از اشتباهات رایج در اسکریپت‌نویسی، عدم تکمیل صحیح دستورات مرکب (\en{Compound Commands}) نظیر \cmd{if} یا \cmd{while} است. برای درک بهتر، پیامد حذف علامت سمی‌کالن (\cmd{;}) را که پس از بخش شرطی در دستور \cmd{if} قرار می‌گیرد، بررسی می‌کنیم:

\begin{latin}
\begin{lstlisting}
#!/bin/bash

# trouble: script to demonstrate common errors

number=1

if [ $number = 1 ] then
    echo "Number is equal to 1."
else
    echo "Number is not equal to 1."
fi
\end{lstlisting}
\end{latin}

خروجی مفسر به شرح زیر خواهد بود:

\begin{latin}
\begin{lstlisting}
[me@linuxbox ~]$ trouble
/home/me/bin/trouble: line 9: syntax error near unexpected token `else'
/home/me/bin/trouble: line 9: `else'
\end{lstlisting}
\end{latin}

در اینجا نیز مشاهده می‌شود که پیام خطا به نقطه‌ای اشاره دارد که پس از رخدادِ خطای واقعی واقع شده است.

تحلیل فرآیند مفسر در این حالت بسیار حائز اهمیت است. همان‌طور که می‌دانیم، ساختار \cmd{if} فهرستی از دستورات را پذیرفته و وضعیت خروج (\en{Exit Status}) آخرین دستور را ارزیابی می‌کند. در اسکریپت فوق، هدف ما این است که فهرست دستورات تنها شامل دستور \cmd{]} (که معادل و مترادف دستور \cmd{test} است) باشد. این دستور، پارامترهای پس از خود را به عنوان فهرستی از آرگومان‌ها دریافت می‌کند که در اینجا شامل چهار مورد است:

\begin{quote}
\lr{\$number} و \lr{=} و \lr{1} و \lr{]}
\end{quote}

با حذف سمی‌کالن، کلمه‌ی کلیدی \cmd{then} به اشتباه به عنوان یکی از آرگومان‌های دستور \cmd{]} در نظر گرفته می‌شود که از نظر نحوی فاقد اشکال است. متعاقباً دستور \cmd{echo} نیز به عنوان دستور بعدی در فهرست دستوراتِ تحت ارزیابیِ \cmd{if} لحاظ می‌گردد. چالش اصلی زمانی بروز می‌کند که مفسر به واژه‌ی \cmd{else} می‌رسد؛ از آنجا که پوسته این واژه را به عنوان یک کلمه رزرو شده (\en{Reserved Word}) و نه یک دستور معمولی می‌شناسد، حضور آن در این موقعیت را غیرمجاز تشخیص داده و پیام خطا صادر می‌کند.

\subsection{خطاهای بسط غیرمنتظره}

برخی خطاها در اسکریپت‌نویسی لینوکس به‌صورت مقطعی رخ می‌دهند؛ به این معنا که اسکریپت در شرایط عادی به‌درستی اجرا می‌شود، اما تحت شرایطی خاص و به دلیل نتایج حاصل از بسط (\en{Expansion})، با شکست مواجه می‌گردد. برای درک این موضوع، سمی‌کالن حذف‌شده در مثال قبل را بازگردانده و مقدار متغیر \cmd{number} را تهی (خالی) فرض می‌کنیم:

\begin{latin}
\begin{lstlisting}
#!/bin/bash

# trouble: script to demonstrate common errors

number=

if [ $number = 1 ]; then
    echo "Number is equal to 1."
else
    echo "Number is not equal to 1."
fi
\end{lstlisting}
\end{latin}

اجرای اسکریپت با این تغییر، خروجی زیر را تولید می‌کند:

\begin{latin}
\begin{lstlisting}
[me@linuxbox ~]$ trouble
/home/me/bin/trouble: line 7: [: =: unary operator expected
Number is not equal to 1.
\end{lstlisting}
\end{latin}

این پیام خطا کمی مبهم به نظر می‌رسد. ریشه‌ی مشکل در نحوه‌ی بسط متغیر \cmd{number} درون دستور \cmd{test} (یا همان \cmd{]}) نهفته است. زمانی که عبارت زیر ارزیابی می‌شود:

\begin{latin}
\begin{lstlisting}
[ $number = 1 ]
\end{lstlisting}
\end{latin}

از آنجا که مقدار متغیر تهی است، پس از مرحله‌ی بسط، دستور عملاً به شکل زیر در می‌آید:

\begin{latin}
\begin{lstlisting}
[  = 1 ]
\end{lstlisting}
\end{latin}

این عبارت از نظر نحوی نامعتبر است. عملگر \cmd{=} یک عملگر دوگانه (\en{Binary Operator}) است و به مقادیری در هر دو طرف خود نیاز دارد؛ اما در اینجا مقدار سمت چپ مفقود است. در چنین حالتی، دستور \cmd{test} به اشتباه تصور می‌کند با یک عملگر یگانی (\en{Unary}) مانند \cmd{-z} روبروست. از آنجا که این ارزیابی با خطا مواجه می‌شود، دستور \cmd{if} وضعیت خروج غیرصفر دریافت کرده و مطابق با آن، بخش \cmd{else} را اجرا می‌کند.

برای رفع این اختلال، کافی است آرگومان‌های دستور \cmd{test} را محصور در دابل‌کوتیشن (\cmd{" "}) قرار دهید:

\begin{latin}
\begin{lstlisting}
[ "$number" = 1 ]
\end{lstlisting}
\end{latin}

در این حالت، پس از بسط، عبارت به شکل زیر در می‌آید که تعداد آرگومان‌های آن صحیح است:

\begin{latin}
\begin{lstlisting}
[ "" = 1 ]
\end{lstlisting}
\end{latin}

علاوه بر متغیرهای تهی، استفاده از کوتیشن برای متغیرهایی که ممکن است حاوی فاصله (\en{Space}) باشند (مانند نام فایل‌ها) نیز حیاتی است.

\begin{note}
\textbf{نکته:} به‌عنوان یک قاعده کلی در توسعه اسکریپت‌های پوسته، همیشه متغیرها و جایگزینی فرمان‌ها (\en{Command Substitution}) را داخل دابل‌کوتیشن قرار دهید، مگر در موارد نادری که صراحتاً به قابلیت «تفکیک کلمات» (\en{Word Splitting}) نیاز داشته باشید.
\end{note}

\section{خطاهای منطقی (Logical Errors)}

برخلاف خطاهای نحوی، خطاهای منطقی مانع از اجرای اسکریپت نمی‌شوند؛ برنامه با موفقیت توسط مفسر پوسته اجرا می‌گردد، اما به دلیل نقص در طراحیِ الگوریتم، خروجی یا نتیجه‌ی مطلوب حاصل نمی‌شود. اگرچه منشأ این خطاها می‌تواند بسیار متنوع باشد، اما در اسکریپت‌نویسی پوسته معمولاً با سه دسته‌ی رایج روبرو هستیم:

\begin{enumerate}
  \item \textbf{عبارات شرطی نادرست:} در بسیاری از موارد، ساختار \cmd{if/then/else} به‌درستی کدگذاری نمی‌شود و در نتیجه منطق برنامه منحرف می‌گردد. این مسئله ممکن است ناشی از معکوس شدن منطقِ شرط یا در نظر نگرفتن تمام حالات ممکن (منطق ناقص) باشد.

  \item \textbf{خطاهای «یک واحد اختلاف» (\en{Off-by-one Errors}):} این خطاها اغلب در حلقه‌هایی رخ می‌دهند که از شمارنده (\en{Counter}) استفاده می‌کنند. برای مثال، برنامه‌نویس ممکن است فراموش کند که شمارش باید از عدد صفر شروع شود یا در نقطه‌ای اشتباه خاتمه یابد. این ناهماهنگی منجر به تکرارِ بیش از حد یا پایان زودهنگام حلقه و از دست رفتن آخرین تکرار می‌شود.

  \item \textbf{موقعیت‌های پیش‌بینی‌نشده:} بخش عمده‌ای از خطاهای منطقی زمانی بروز می‌کنند که برنامه با داده‌ها یا شرایطی مواجه می‌شود که توسعه‌دهنده از قبل برای آن‌ها تدبیری نیندیشیده است. همان‌طور که پیش‌تر اشاره شد، بسط‌های غیرمنتظره (مانند نام فایلی که حاوی کاراکتر فاصله است و به اشتباه به چندین آرگومان مجزا تبدیل می‌شود) نمونه‌ای بارز از این وضعیت‌های پیش‌بینی‌نشده در دنیای لینوکس هستند.
\end{enumerate}

\section{برنامه‌نویسی تدافعی (Defensive Programming)}

در توسعه‌ی اسکریپت، اعتبارسنجی فرضیات نقشی حیاتی ایفا می‌کند. این رویکرد به معنای ارزیابی دقیق وضعیت خروج (\en{Exit Status}) تمامی دستورات و برنامه‌های به‌کاررفته است. برای درک اهمیت این موضوع، یک نمونه‌ی واقعی را بررسی می‌کنیم: مدیر سیستمی برای انجام عملیات نگهداری روی یک سرور حساس، اسکریپتی شامل دو خط زیر نوشت:

\begin{latin}
\begin{lstlisting}
cd $dir_name
rm *
\end{lstlisting}
\end{latin}

در نگاه اول، این دو خط فاقد ایراد فنی هستند، اما تنها به شرطی که دایرکتوریِ تعریف‌شده در متغیر \cmd{dir\_name} واقعاً وجود داشته باشد. اگر این مسیر وجود نداشته باشد، دستور \cmd{cd} با خطا مواجه شده و مفسر به خط بعد می‌رود؛ در نتیجه، تمام فایل‌های موجود در دایرکتوری کاری فعلی (\en{Current Working Directory}) حذف می‌شوند. این خطای منطقی می‌تواند منجر به نابودی بخش‌های حیاتی سیستم شود.

برای ارتقای امنیت و پایداری این طراحی، می‌توان چندین گام برداشت. نخست، با استفاده از دابل‌کوتیشن، از بسط صحیح متغیر مطمئن می‌شویم و اجرای دستور \cmd{rm} را به موفقیتِ \cmd{cd} منوط می‌کنیم:

\begin{latin}
\begin{lstlisting}
cd "$dir_name" && rm *
\end{lstlisting}
\end{latin}

با این تغییر، اگر \cmd{cd} شکست بخورد، \cmd{rm} اجرا نخواهد شد. با این حال، اگر متغیر \cmd{dir\_name} تعریف نشده یا خالی باشد، دستور \cmd{cd} ممکن است کاربر را به دایرکتوری خانگی (\en{Home}) هدایت کند. برای جلوگیری از این سناریو، باید ابتدا وجود دایرکتوری را احراز کنیم:

\begin{latin}
\begin{lstlisting}
[[ -d "$dir_name" ]] && cd "$dir_name" && rm *
\end{lstlisting}
\end{latin}

در اسکریپت‌های حرفه‌ای، بهترین رویکرد این است که منطقی برای توقف برنامه و گزارش خطا در صورت بروز شرایط پیش‌بینی‌نشده اضافه شود:

\begin{latin}
\begin{lstlisting}
# Delete files in directory $dir_name
if [[ ! -d "$dir_name" ]]; then
    echo "No such directory: '$dir_name'" >&2
    exit 1
fi

if ! cd "$dir_name"; then
    echo "Cannot cd to '$dir_name'" >&2
    exit 1
fi

if ! rm *; then
    echo "File deletion failed. Check results" >&2
    exit 1
fi
\end{lstlisting}
\end{latin}

در این ساختار، ابتدا موجودیت مسیر بررسی شده و سپس موفقیتِ تغییر دایرکتوری کنترل می‌گردد. در صورت بروز هرگونه اختلال، یک پیام توضیحی به خروجی خطای استاندارد (\en{stderr}) ارسال شده و اسکریپت با وضعیت خروج \cmd{1} (نشانه شکست) خاتمه می‌یابد.

\section{مکانیزم‌های توقف خودکار: set -e، set -u و set -o PIPEFAIL}

یکی از ویژگی‌های پیش‌فرض در پوسته \cmd{bash} این است که اگر دستوری در حین اجرای اسکریپت با شکست مواجه شود (به‌جز خطاهای نحوی در خودِ اسکریپت)، مفسر بدون توقف به اجرای دستور بعدی ادامه می‌دهد. در بسیاری از موارد، این رفتار نامطلوب است؛ به همین دلیل استاندارد \en{POSIX} و به‌تبع آن \cmd{bash} راهکارهایی را برای مدیریت این وضعیت ارائه داده‌اند. پوسته \cmd{bash} از تنظیماتی بهره می‌برد که امکان مدیریت خودکار خطاها را فراهم می‌کند؛ به این معنا که با فعال‌سازی آن‌ها، اسکریپت به محض بازگشت وضعیت خروج (\en{Exit Code}) غیرصفر از هر دستور (با رعایت برخی استثنائات منطقی)، متوقف می‌شود.

برای فعال‌سازی این قابلیت، معمولاً دستور \cmd{set -e} در ابتدای اسکریپت قرار می‌گیرد. علاوه بر این، بسیاری از استانداردهای نوین اسکریپت‌نویسی بر استفاده هم‌زمان از این ویژگی و چند تنظیم مرتبط دیگر تأکید دارند؛ از جمله \cmd{set -u} که در صورت ارجاع به متغیرهای تعریف‌نشده (\en{Unset}) باعث توقف اسکریپت می‌شود، و \cmd{set -o PIPEFAIL} که موجب می‌گردد در صورت شکست هر یک از اجزای یک پایپ‌لاین (\en{Pipeline})، کل فرآیند متوقف شود (برخلاف حالت پیش‌فرض که فقط وضعیت خروج آخرین دستورِ پایپ ملاک است).

با این وجود، تکیه مطلق بر این ویژگی‌ها توصیه نمی‌شود. رویکرد اصولی آن است که مدیریت خطا به‌صورت آگاهانه طراحی شود و نباید از \cmd{set -e} به‌عنوان جایگزینی برای الگوهای صحیح برنامه‌نویسی استفاده کرد.

در بخش پرسش‌های متداول (\en{FAQ}) شماره ۱۰۵ مستندات \en{Bash}، تحلیل زیر در این باره ارائه شده است:

\begin{latin}
\url{http://mywiki.wooledge.org/BashFAQ/105}
\end{latin}

\begin{quote}
«قابلیت \cmd{set -e} تلاشی برای افزودن مکانیزم «کشف خودکار خطا» به پوسته بود. هدف اصلی این بود که پوسته با بروز هر خطایی متوقف شود تا برنامه‌نویس ناچار نباشد پس از هر دستور حیاتی، عبارت \cmd{|| exit 1} را به‌صورت دستی اضافه کند.»
\end{quote}

تحقق این هدف بسیار پیچیده است؛ چرا که بسیاری از دستورات به‌طور عامدانه وضعیت خروج غیرصفر برمی‌گردانند. برای مثال در عبارت:

\begin{latin}
\begin{lstlisting}
if [ -d /foo ]; then ...; else ...; fi
\end{lstlisting}
\end{latin}

بدیهی است که نمی‌خواهیم در صورت عدم وجود دایرکتوری و بازگشت وضعیت خروج غیرصفر توسط دستور \cmd{[ -d /foo ]}، اسکریپت متوقف شود؛ چرا که قصد داریم این وضعیت را در بلوک \cmd{else} مدیریت کنیم.

به همین سبب، توسعه‌دهندگان ناگزیر به وضع مجموعه‌ای از قوانین استثنایی شدند؛ قوانینی نظیر «دستوراتی که بخشی از شرط \cmd{if} هستند، ایمن تلقی می‌شوند» یا «دستورات میانی در یک پایپ‌لاین (\en{Pipeline})، به‌جز آخرین دستور، مشمول توقف نمی‌شوند.»

این قوانین بسیار پیچیده هستند و حتی در پوشش برخی سناریوهای ساده نیز ناتوان می‌مانند. چالش بزرگ‌تر اینجاست که این ضوابط در نسخه‌های مختلف \en{Bash} تغییر می‌کنند، زیرا این پوسته همواره در تلاش است تا تعاریف منعطف و گاه مبهم \en{POSIX} از این ویژگی را دنبال کند. این عدم قطعیت هنگام درگیر شدن زیرپوسته‌ها (\en{Subshells}) دوچندان می‌شود؛ در این حالت، رفتار سیستم بسته به اینکه \en{Bash} در حالت انطباق با \en{POSIX} فراخوانی شده باشد یا خیر، متفاوت خواهد بود.

\section{ابزار عیب‌یابی ShellCheck}

برنامه‌ای کاربردی به نام \cmd{ShellCheck} در مخازن رسمی اکثر توزیع‌های لینوکس موجود است که وظیفه‌ی تحلیل ایستای اسکریپت‌های پوسته را بر عهده دارد. این ابزار قادر است طیف گسترده‌ای از خطاهای نحوی، منطقی و شیوه‌های نادرست اسکریپت‌نویسی (\en{Anti-patterns}) را شناسایی کند. بهره‌گیری از این ابزار برای تضمین کیفیت و پایداری کدهای توسعه‌یافته بسیار توصیه می‌شود.

برای بررسی اسکریپتی که دارای خط تفسیر (\en{Shebang}) است، به‌سادگی دستور زیر را اجرا کنید:

\begin{latin}
\begin{lstlisting}
shellcheck my_script
\end{lstlisting}
\end{latin}

\cmd{ShellCheck} به‌طور خودکار نوع پوسته (\en{Shell}) را بر اساس محتوای \en{shebang} تشخیص می‌دهد. جهت تحلیل فایل‌هایی که فاقد این خط هستند (مانند کتابخانه‌های توابع)، می‌توان از پارامتر \cmd{-s} برای تعیین صریح نوع پوسته استفاده کرد:

\begin{latin}
\begin{lstlisting}
shellcheck -s bash my_library
\end{lstlisting}
\end{latin}

در اینجا گزینه \cmd{-s} مخفف \en{shell} بوده و برای مشخص کردن مفسر مورد نظر (مانند \cmd{bash}، \cmd{sh} یا \cmd{ksh}) به کار می‌رود. برای کسب اطلاعات بیشتر و استفاده از نسخه‌ی آنلاین این ابزار، می‌توانید به وب‌سایت رسمی آن در نشانی \en{shellcheck.net} مراجعه کنید.

\section{مخاطرات نام‌گذاری فایل‌ها}

اختلال دیگری که در اسکریپت‌های حذف فایل وجود دارد، کمتر به چشم می‌آید اما می‌تواند پیامدهای بسیار خطرناکی داشته باشد. از دیدگاه بسیاری از متخصصان، سیستم‌های یونیکس و شبه‌یونیکس دارای یک چالش طراحی بنیادین در حوزه نام‌گذاری فایل‌ها هستند. یونیکس در این زمینه بسیار منعطف عمل می‌کند؛ به طوری که عملاً تنها دو نویسه (\en{Character}) مجاز به حضور در نام فایل نیستند: نخست کاراکتر \cmd{/} که برای تفکیک اجزای مسیر (\en{Path}) استفاده می‌شود و دوم نویسه‌ی \en{null} که برای نشانه‌گذاری پایان رشته‌ها کاربرد دارد.

سایر نویسه‌ها کاملاً مجاز هستند؛ از جمله فاصله‌ها (\en{Spaces})، تب‌ها (\en{Tabs})، خطِ جدید (\en{Newline})، خط‌تیره در ابتدای نام، بازگشت به ابتدای سطر (\en{Carriage Return}) و غیره.

نکته‌ی بحرانی، وجود خط‌تیره (\cmd{-}) در ابتدای نام فایل است. برای مثال، ایجاد فایلی با نام \file{-rf~} در یونیکس کاملاً مجاز است. لحظه‌ای تصور کنید اگر این نام فایل به‌عنوان ورودی به دستور \cmd{rm} ارسال شود، چه اتفاقی رخ می‌دهد؛ مفسر دستور، نام فایل را به‌عنوان گزینه‌های عملیاتی (\cmd{-r} برای حذف بازگشتی و \cmd{-f} برای اجبار) و علامت \cmd{\textasciitilde} را به‌عنوان دایرکتوری خانگی کاربر تفسیر می‌کند.

برای ایمن‌سازی اسکریپت در برابر این آسیب‌پذیری، باید دستور \cmd{rm} را از حالت ساده به شکل زیر تغییر دهیم:

\begin{latin}
\begin{lstlisting}
rm ./*
\end{lstlisting}
\end{latin}

افزودن پیشوند \cmd{./} باعث می‌شود تا نام فایل‌هایی که با خط‌تیره شروع می‌شوند، دیگر به‌عنوان «گزینه» (\en{Option}) توسط دستور تفسیر نشوند، چرا که اکنون مسیر با یک نقطه شروع شده است. به‌طور کلی، توصیه می‌شود هنگام استفاده از نویسه‌های جانشین یا همان \en{Wildcard}ها (نظیر \cmd{*} و \cmd{?})، همیشه از پیشوند \cmd{./} استفاده کنید تا از تفسیر نادرست آن‌ها توسط دستورات جلوگیری شود. این قاعده شامل موارد مشابهی مانند \cmd{*.pdf} یا \cmd{*.mp3} نیز می‌گردد.

\section{نام‌گذاری فایل‌های قابل حمل (Portable Filenames)}

برای تضمین سازگاری نام فایل‌ها میان پلتفرم‌های گوناگون (مانند توزیع‌های مختلف لینوکس، یونیکس و سایر سیستم‌عامل‌ها)، رعایت احتیاط در انتخاب نویسه‌ها امری حیاتی است. استانداردی تحت عنوان \en{POSIX Portable Filename Character Set} تدوین شده است که با پیروی از آن، احتمال عملکرد صحیح نام فایل در سیستم‌های مختلف به حداکثر می‌رسد.

این استاندارد بسیار ساده و محدود است؛ بر این اساس، تنها نویسه‌های مجاز برای استفاده عبارتند از: حروف بزرگ (\en{A-Z})، حروف کوچک (\en{a-z})، اعداد (\en{0-9})، نقطه (\cmd{.})، خط‌تیره (\cmd{-}) و زیرخط (\cmd{\_}). همچنین، این استاندارد مؤکداً پیشنهاد می‌کند که نام فایل‌ها هرگز با کاراکتر خط‌تیره شروع نشوند تا از بروز تداخل‌های نحوی در محیط پوسته جلوگیری شود.

\section{اعتبارسنجی ورودی‌ها}

یک اصل بنیادین در توسعه نرم‌افزار حرفه‌ای حکم می‌کند که اگر برنامه‌ای ورودی دریافت می‌کند، باید توانایی مدیریت هر نوع داده‌ی ورودی را داشته باشد. این امر مستلزم بررسی و پالایش دقیق داده‌ها است تا اطمینان حاصل شود که تنها مقادیر معتبر به مراحل بعدی پردازش راه می‌یابند. در فصل پیشین و هنگام بررسی دستور \cmd{read}، نمونه‌ای از این رویکرد را در قالب شرط زیر برای اعتبارسنجی گزینه‌های منو مشاهده کردیم:

\begin{latin}
\begin{lstlisting}
[[ $REPLY =~ ^[0-3]$ ]]
\end{lstlisting}
\end{latin}

این عبارتِ شرطی بسیار دقیق عمل می‌کند و تنها زمانی وضعیت خروج صفر را برمی‌گرداند که رشته‌ی ورودی کاربر دقیقاً یک عدد در بازه صفر تا سه باشد؛ در غیر این صورت، هر ورودی دیگری مردود اعلام می‌شود. نگارش چنین شرط‌های سخت‌گیرانه‌ای اگرچه ممکن است چالش‌برانگیز باشد، اما برای تولید یک اسکریپت با کیفیت و پایدار (\en{Robust}) ضرورتی انکارناپذیر است.

\section{طراحی اسکریپت و محدودیت زمان}

در دنیای طراحی صنعتی، اصلی وجود دارد که می‌گوید عمق و کیفیت طراحی یک پروژه، تابعی از زمانِ اختصاص‌یافته به طراح است. اگر برای طراحی ابزاری جهت «دفع حشرات» تنها پنج دقیقه فرصت داشته باشید، نتیجه‌ی کار یک مگس‌کش ساده خواهد بود؛ اما با اختصاص پنج ماه زمان، ممکن است یک «سیستم لیزری هوشمند» ابداع کنید.

این منطق در دنیای برنامه‌نویسی و اسکریپت‌نویسی پوسته نیز کاملاً صادق است. گاهی یک اسکریپت موقت و ساده (\en{Quick and Dirty}) که تنها یک‌بار توسط خودِ برنامه‌نویس استفاده می‌شود، کفایت می‌کند. چنین ابزارهایی باید به سرعت و با کمترین هزینه نوشته شوند و نیازی به مستندسازی گسترده یا مکانیزم‌های دفاعی پیچیده ندارند. اما اگر قرار است اسکریپتی در محیط‌های عملیاتی (\en{Production}) مستقر شود و به‌طور مکرر توسط کاربران متعدد مورد استفاده قرار گیرد، فرآیند توسعه‌ی آن باید بسیار دقیق، نظام‌مند و با رعایت تمامی استانداردهای ایمنی انجام پذیرد.

\section{متدولوژی آزمایش (Testing)}

آزمایش و اعتبارسنجی، گامی حیاتی در فرآیند توسعه‌ی هر نرم‌افزار، از جمله اسکریپت‌نویسی پوسته محسوب می‌شود. در زیست‌بوم متن‌باز (\en{Open Source})، شعار معروف «زود منتشر کن، زیاد منتشر کن» (\en{Release Early, Release Often}) بر همین منطق استوار است. با انتشار زودهنگام و مستمر، نرم‌افزار در معرض استفاده و ارزیابی گسترده‌تری قرار می‌گیرد؛ چرا که تجربه نشان داده است شناسایی و رفع باگ‌ها در مراحل اولیه‌ی توسعه، به‌مراتب ساده‌تر و کم‌هزینه‌تر از مراحل نهایی است.

در فصل ۲۶، آموختیم که چگونه می‌توان از \en{Stub}ها برای تأیید جریان اجرای برنامه استفاده کرد. این تکنیک از همان نخستین مراحل کدنویسی، ابزاری ارزشمند برای پایش پیشرفت پروژه و اطمینان از صحت منطق اسکریپت است.

بیایید چالش حذف فایل را که پیش‌تر بررسی کردیم، مجدداً در نظر بگیریم و روشی ایمن برای آزمایش آن طراحی کنیم. تست کردن مستقیمِ قطعه‌کد اصلی خطرناک است، زیرا هدف آن حذف واقعی داده‌هاست؛ اما می‌توان با اعمال تغییراتی، فرآیند شرط را بدون ریسک انجام داد:

\begin{latin}
\begin{lstlisting}
if [[ -d $dir_name ]]; then
    if cd $dir_name; then
        echo rm * # TESTING
    else
        echo "cannot cd to '$dir_name'" >&2
        exit 1
    fi
else
    echo "no such directory: '$dir_name'" >&2
    exit 1
fi
exit # TESTING
\end{lstlisting}
\end{latin}

در این ساختار، از آنجا که بخش‌های مدیریت خطا پیام‌های خروجی گویایی صادر می‌کنند، نیازی به افزودن پیام‌های جدید نیست. تغییر کلیدی، قرار دادن دستور \cmd{echo} پیش از دستور \cmd{rm} است. این کار باعث می‌شود به جای اجرای واقعیِ دستور حذف، خودِ دستور به همراه لیست پارامترهای بسط‌یافته‌ی آن (\en{Expansion}) در خروجی نمایش داده شود. این رویکرد امکان بررسی دقیق رفتار اسکریپت را در محیطی ایزوله فراهم می‌کند. همچنین در انتهای قطعه‌کد، دستور \cmd{exit} اضافه شده است تا از اجرای ادامه‌ی اسکریپت در طول فرآیند تست جلوگیری شود؛ هرچند ضرورت این کار به ساختار کلی اسکریپت بستگی دارد.

علاوه بر این، از کامنت‌هایی با محتوای \cmd{\# TESTING} استفاده شده است. این برچسب‌ها به توسعه‌دهنده کمک می‌کنند تا پس از اتمام مراحل عیب‌یابی، به‌سرعت تغییرات موقت را شناسایی و پیش از استقرار نهایی اسکریپت، آن‌ها را حذف کند.

\subsection{طراحی سناریوهای آزمایش (Test Cases)}

برای انجام یک ارزیابی اثربخش، تدوین و اجرای موارد آزمایشی مناسب ضرورتی کلیدی دارد. این فرآیند از طریق انتخاب دقیق داده‌های ورودی یا بازسازی شرایط عملیاتی خاصی صورت می‌پذیرد که «حالت‌های مرزی» (\en{Edge Cases}) و سناریوهای استثنایی را هدف قرار می‌دهند. در قطعه‌کد مورد بحث، ما باید رفتار اسکریپت را در سه وضعیت عملیاتی مجزا تحلیل کنیم:

\begin{enumerate}
  \item متغیر \cmd{dir\_name} حاوی مسیر یک دایرکتوری موجود و معتبر باشد.
  \item متغیر \cmd{dir\_name} به مسیری اشاره کند که در سیستم وجود ندارد.
  \item متغیر \cmd{dir\_name} کاملاً تهی (\en{Empty}) باشد.
\end{enumerate}

با اجرای اسکریپت در هر یک از این شرایط، «پوشش آزمایش» (\en{Test Coverage}) مناسبی ایجاد می‌شود که ریسک رفتارهای پیش‌بینی‌نشده را به حداقل می‌رساند.

همانند فرآیند طراحی، آزمایش نیز تابعی از محدودیت‌های زمانی است. لزومی ندارد تمامی جزئیات و عملکردهای فرعی یک اسکریپت تحت شرط‌های گسترده قرار گیرند؛ بلکه هنر اصلی در «اولویت‌بندی» (\en{Prioritization}) است. از آنجا که نقص در عملکرد این قطعه‌کد خاص می‌تواند منجر به تخریب جبران‌ناپذیر داده‌ها شود، سرمایه‌گذاری زمانی برای طراحی دقیق و آزمایش وسواس‌گونه‌ی آن کاملاً توجیه‌پذیر و ضروری است.

\section{متدولوژی عیب‌یابی (Debugging)}

چنانچه نتایج حاصل از آزمایش‌ها بیانگر وجود اختلال در عملکرد اسکریپت باشد، گام بعدی ورود به فرآیند عیب‌یابی است. بروز «مشکل» بدین معناست که اسکریپت برخلاف انتظارات و منطق طراحی‌شده توسط برنامه‌نویس رفتار می‌کند. در چنین شرایطی، باید با رویکردی تحلیل‌گرانه مشخص کنیم که اسکریپت در عمل چه توالی عملیاتی را طی می‌کند و منشأ این انحراف کجاست. شناسایی ریشه‌ی باگ‌ها گاه مستلزم انجام تحقیقات دقیق و ایفای نقش یک کارآگاه فنی است.

یک اسکریپت با طراحی استاندارد، خود ابزاری برای تسهیل این فرآیند است. برنامه‌نویسی با رویکرد تدافعی باعث می‌شود تا اسکریپت شرایط غیرعادی را به‌سرعت تشخیص داده و بازخوردهای (\en{Feedback}) ارزشمندی را به کاربر یا توسعه‌دهنده ارائه دهد. با این حال، در مواجهه با مشکلات پیچیده و رفتارهای غیرمنتظره، تکیه بر پیام‌های خطای ساده کافی نبوده و استفاده از تکنیک‌های پیشرفته‌تر عیب‌یابی اجتناب‌ناپذیر می‌گردد.

\subsection{ایزوله‌سازی محدوده خطا}

در مواجهه با اسکریپت‌های پیچیده و طولانی، جداسازی بخشی از کد که منشأ بروز رفتار غیرعادی است، راهکاری بسیار کارآمد محسوب می‌شود. لزوماً بخش ایزوله‌شده محل دقیق وقوع خطای اصلی نیست، اما این فرآیند دیدگاه‌های ارزشمندی درباره‌ی ریشه‌ی واقعی مشکل در اختیار برنامه‌نویس قرار می‌دهد. یکی از متداول‌ترین روش‌ها برای این کار، تبدیل بخش‌هایی از کد به «توضیحات» یا همان کامنت‌گذاری (\en{Commenting Out}) است.

به‌عنوان مثال، در قطعه‌کد حذف فایل، می‌توان با غیرفعال کردن موقت بخش‌های منطقی، تأثیر آن‌ها بر رفتار کلی اسکریپت و خطاهای گزارش‌شده بررسی کرد:

\begin{latin}
\begin{lstlisting}
if [[ -d $dir_name ]]; then
    if cd $dir_name; then
        rm *
    else
        echo "cannot cd to '$dir_name'" >&2
        exit 1
    fi
# else
#   echo "no such directory: '$dir_name'" >&2
#   exit 1
fi
\end{lstlisting}
\end{latin}

با قرار دادن نویسه‌ی \cmd{\#} در ابتدای خطوط، از اجرای آن بخش منطقی جلوگیری می‌شود. پس از این تغییر، می‌توان آزمون را مجدداً تکرار کرد تا مشخص شود که آیا حذفِ موقت آن بخش از کد، تأثیری بر رفع یا تغییر ماهیت خطا داشته است یا خیر. این رویکرد به شما اجازه می‌دهد تا دایره‌ی جستجو را برای یافتن باگ‌های منطقی تنگ‌تر کنید.

\subsection{تکنیک ردیابی عملیات}

بسیاری از خطاها ناشی از انحراف در جریان منطقی اسکریپت هستند؛ به این معنا که بخش‌هایی از کد یا هرگز فراخوانی نمی‌شوند، یا با ترتیب و زمان‌بندی نادرست به اجرا درمی‌آیند. برای درک دقیق توالیِ واقعیِ اجرای دستورات، از تکنیکی تحت عنوان ردیابی (\en{Tracing}) استفاده می‌کنیم.

یکی از متداول‌ترین روش‌های ردیابی، درج پیام‌های اطلاع‌رسان در نقاط کلیدی اسکریپت است تا موقعیت لحظه‌ایِ اجرای کد مشخص شود. برای نمونه، می‌توانیم پیام‌هایی را به قطعه‌کد خود اضافه کنیم:

\begin{latin}
\begin{lstlisting}
echo "preparing to delete files" >&2
if [[ -d $dir_name ]]; then
    if cd $dir_name; then
echo "deleting files" >&2
        rm *
    else
        echo "cannot cd to '$dir_name'" >&2
        exit 1
    fi
else
    echo "no such directory: '$dir_name'" >&2
    exit 1
fi
echo "file deletion complete" >&2
\end{lstlisting}
\end{latin}

در این رویکرد، پیام‌های ردیابی به خروجی خطای استاندارد (\en{stderr}) ارسال می‌شوند تا با خروجی معمول اسکریپت (\en{stdout}) تداخل پیدا نکنند. همچنین، برای سهولت در شناسایی و حذفِ این خطوط پس از اتمام فرآیند عیب‌یابی، آن‌ها را بدون تورفتگی و در ابتدای سطر درج می‌کنیم تا از بدنه‌ی اصلی کد متمایز باشند.

با اجرای اسکریپت در حالت ردیابی، اکنون می‌توان توالی دقیق عملیات را مشاهده کرد:

\begin{latin}
\begin{lstlisting}
[me@linuxbox ~]$ deletion-script
preparing to delete files
deleting files
file deletion complete
[me@linuxbox ~]$
\end{lstlisting}
\end{latin}

پوسته \cmd{bash} علاوه بر روش‌های دستی، مکانیزم ردیابی داخلی قدرتمندی دارد که از طریق گزینه \cmd{-x} (چه در خط فرمان و چه با دستور \cmd{set}) فعال می‌شود. برای فعال‌سازی ردیابی در کل یک اسکریپت، کافی است گزینه \cmd{-x} را به خط شِبنگ (\en{Shebang}) اضافه کنید. به نمونه زیر توجه کنید:

\begin{latin}
\begin{lstlisting}
#!/bin/bash -x

# trouble: script to demonstrate common errors

number=1

if [ $number = 1 ]; then
    echo "Number is equal to 1."
else
    echo "Number is not equal to 1."
fi
\end{lstlisting}
\end{latin}

پس از اجرا، خروجی به صورت زیر نمایش داده می‌شود:

\begin{latin}
\begin{lstlisting}
[me@linuxbox ~]$ trouble
+ number=1
+ '[' 1 = 1 ']'
+ echo 'Number is equal to 1.'
Number is equal to 1.
\end{lstlisting}
\end{latin}

در این حالت، تمامی دستورات پس از انجام عملیات بسط (\en{Expansion}) نمایش داده می‌شوند. نویسه‌ی \cmd{+} در ابتدای خطوط، خروجیِ ردیابی را از خروجی استاندارد برنامه متمایز می‌کند. این کاراکتر، مقدار پیش‌فرض متغیر سیستمی \cmd{PS4} (\en{Prompt String 4}) است. شما می‌توانید برای افزایش کارایی فرآیند عیب‌یابی، محتوای این متغیر را شخصی‌سازی کنید؛ برای مثال، آن را طوری تنظیم کنید که «شماره خط» اسکریپت را در کنار هر دستور نمایش دهد.

توجه داشته باشید که برای جلوگیری از بسط زودهنگام (\en{Early Expansion}) متغیر، حتماً باید از تک‌کوتیشن (\en{Single Quote}) استفاده کنید:

\begin{latin}
\begin{lstlisting}
[me@linuxbox ~]$ export PS4='$LINENO + '
[me@linuxbox ~]$ trouble
5 + number=1
7 + '[' 1 = 1 ']'
8 + echo 'Number is equal to 1.'
Number is equal to 1.
\end{lstlisting}
\end{latin}

جهت اجرای عملیات ردیابی (\en{Tracing}) بر روی یک قطعه‌کد خاص به جای کل اسکریپت، می‌توان از دستور داخلی \cmd{set} همراه با گزینه \cmd{-x} استفاده کرد. این رویکرد به شما اجازه می‌دهد تا فرآیند عیب‌یابی را دقیقاً بر روی ناحیه مشکوک متمرکز کنید:

\begin{latin}
\begin{lstlisting}
#!/bin/bash

# trouble: script to demonstrate common errors

number=1

set -x # Turn on tracing
if [ $number = 1 ]; then
    echo "Number is equal to 1."
else
    echo "Number is not equal to 1."
fi
set +x # Turn off tracing
\end{lstlisting}
\end{latin}

در این متدولوژی، از \cmd{set -x} برای شروع ردیابی و از \cmd{set +x} برای خاتمه دادن به آن استفاده می‌شود. بهره‌گیری از این تکنیک برای تحلیل بخش‌های مجزای یک اسکریپت معیوب، سرعت شناسایی خطاهای منطقی را به‌شدت افزایش می‌دهد.

\subsection{پایش متغیرها در زمان اجرا}

در کنار فرآیند ردیابی، نمایش محتوای متغیرها ابزاری بسیار کارآمد برای تحلیل مکانیسم‌های داخلی اسکریپت در حین اجراست. با درج هوشمندانه‌ی دستورهای \cmd{echo} در نقاط حساس، می‌توان وضعیت لحظه‌ای داده‌ها را مشاهده کرد.

\begin{latin}
\begin{lstlisting}
#!/bin/bash

# trouble: script to demonstrate common errors

number=1

echo "number=$number" # DEBUG
set -x # Turn on tracing
if [ $number = 1 ]; then
    echo "Number is equal to 1."
else
    echo "Number is not equal to 1."
fi
set +x # Turn off tracing
\end{lstlisting}
\end{latin}

در این نمونه‌ی ساده، مقدار متغیر \cmd{number} را در خروجی چاپ کرده و سطر افزوده شده را با کامنت \cmd{\# DEBUG} نشانه‌گذاری کرده‌ایم تا در مراحل نهایی پاکسازی، به‌راحتی قابل شناسایی و حذف باشد. این تکنیک به‌ویژه هنگام تحلیل رفتار حلقه‌ها (\en{Loops}) و بررسی صحت عملیات محاسباتی در اسکریپت‌های پیچیده، نقشی کلیدی در شناسایی انحرافات منطقی ایفا می‌کند.

\section{جمع‌بندی}

در این فصل، تنها بخشی از چالش‌هایی را که ممکن است در طول توسعه‌ی اسکریپت رخ دهند، بررسی کردیم. بی‌تردید موانع پیش‌روی یک توسعه‌دهنده بسیار فراتر از این موارد است؛ اما متدولوژی‌های ارائه‌شده، بستر لازم برای شناسایی و رفع اکثر خطاهای رایج را فراهم می‌سازند. عیب‌یابی (\en{Debugging}) هنری ظریف است که با ممارست و تجربه صیقل می‌یابد؛ این مهارت هم شامل پیشگیری از بروز خطا (از طریق آزمایش‌های مستمر در حین توسعه) و هم توانمندی در ردیابی و ریشه‌یابی اختلالات (با بهره‌گیری مؤثر از تکنیک‌های \en{Tracing}) می‌گردد.

\section{منابع مطالعه تکمیلی}

ویکی‌پدیا مقالات کوتاهی درباره خطاهای نحوی (\en{syntactic errors}) و خطاهای منطقی (\en{logical errors}) دارد:

\begin{latin}
\url{http://en.wikipedia.org/wiki/Syntax_error}\\
\url{http://en.wikipedia.org/wiki/Logic_error}
\end{latin}

منابع آنلاین متعددی برای بررسی جنبه‌های فنی برنامه‌نویسی \en{Bash} وجود دارد:

\begin{latin}
\url{http://mywiki.wooledge.org/BashPitfalls}\\
\url{http://tldp.org/LDP/abs/html/gotchas.html}\\
\url{http://www.gnu.org/software/bash/manual/html_node/Reserved-Word-Index.html}
\end{latin}

دیوید ویلر (\en{David Wheeler}) بحث بسیار مفیدی درباره چالش نام‌گذاری فایل‌ها در یونیکس و نحوه کدنویسی اسکریپت‌های پوسته (\en{shell scripts}) برای مدیریت آن‌ها ارائه کرده است:

\begin{latin}
\url{https://www.dwheeler.com/essays/filenames-in-shell.html}
\end{latin}

برای اشکال‌زدایی (\en{debugging}) در سطوح پیشرفته، می‌توانید از ابزار اشکال‌زدای مفسر \en{Bash} (\en{Bash Debugger}) استفاده کنید:

\begin{latin}
\url{http://bashdb.sourceforge.net/}
\end{latin}